\section{Discussion}\label{sec:discussion}

\subsection{Why the results are what they are}\label{sec:disc-why}

The shape of the front is set by one fact that Proposition~\ref{prop:shift} and Figure~\ref{fig:shift} expose: the loss can ask for as much light as it likes, and the hologram can deliver only \eMaxRatioPct{} percent more than the start (\eMaxRatio{} times the efficiency of the start at the largest weight). Every extra percent of efficiency is paid for by a large drop of uniformity, so the front is steep and the useful part of the trade-off, where the uniformity is still at least 0.9, spans an increase of the efficiency of
only \eGainUNinePp{} percentage points. This is why the single points of weighted GS and phase-fixed WGS, which sit at the uniform end, are already close to the best that the task allows, and why the hypervolume with the reference $(0,0)$ is dominated by that end: a point $(e,u)$ near $(0.8,1)$ covers $0.8$ of an area whose total is one. The hypervolume of the front exceeds that of the union of the published points by \hvDiff{} not because the uniform end improved but because the front adds points of
higher efficiency at intermediate uniformity, where no published point lies. The sensitivity analysis with the references $(0.5,0)$ and $(0.7,0)$ shows that the advantage is robust (\hvRefHalfRel{} and \hvRefSevenRel{} percent of the union against \hvRatioPct{} percent), and that it is larger in relative terms the more the reference discounts the uniform end. A reader who cares about the uniform end only should read the dominance counts and not the hypervolume.

The front also depends on its start. The ablation shows that the start from phase-fixed WGS is better than a start from plain GS or from the random superposition, and the diagnostic shows that it is not the best possible: weighted GS after 100 iterations reaches a hologram with more light in the traps at nearly the same flatness, and the loss itself is lower for that hologram than for the front at \diagGswHundredMin{} to \diagGswHundredMax{} of the weights. The conjugate-gradient iteration with 15 steps per weight does not leave the basin of its start, and the first part of the front
is that of the basin of phase-fixed WGS, whose efficiency is lower than that of the converged iterates of weighted GS. When the front starts from the converged hologram no front point is dominated and the hypervolume is unchanged within the interval. The loss is not convex, the stationarity of Proposition~\ref{prop:mono}(b) is a necessary condition only, and the engine's front is a set of local solutions: the warm start keeps them on one branch, which the monotone counts of Section~\ref{sec:res-front} support, and does not make the branch the best one.

The ablation reads consistently with the theory. The cold start loses because the warm start's initial gradient is proportional to the step in the weight (Proposition~\ref{prop:warm}); the unscaled loss loses because it asks for a raise of the traps that is $M/e_0=\rawFactor{}$ times larger at the same weight, which Proposition~\ref{prop:shift} makes explicit through the constant $\delta_\lambda$, proportional to $\sum T^2/e_0$; and the plateau in the steps per weight and the number of weights above 20 follows from a front that is smooth.

\subsection{Failure cases and negative results}\label{sec:disc-failures}

Four results of the paper are negative or partial and are reported as such. First, the front does not dominate the point of weighted GS or of phase-fixed WGS in most seeds: the fourth hypothesis is not confirmed, and on this task the continuation adds the efficient side of the trade-off and not a better uniform end. Second, \frontDomMin{} to \frontDomMax{} of the \nLambda{} points of the front are dominated by iterates that the published methods reach, in particular by weighted GS after 100 iterations,
so the set that the front draws is an inner approximation of the attainable set, not the Pareto front of the problem; the converged start removes the domination with no change in area that five seeds can resolve, and the module as it stands does not use it. Third, the stored fronts of the benchmark meet the published union on \benchPassHundred{} of the \benchNHundred{} geometries at 100 iterations (the margins run from \benchWorstHundred{} to \benchBestHundred{}) and on \benchPassThirty{} of \benchNThirty{} at 30 iterations; the stored fronts end at weights well below one (the largest stored weight of each geometry is in Table~\ref{tab:bench}), which is consistent with the 20 seconds that the benchmark's converged row allows for 300 iterations per weight, they hold between one and seven points, and the published medians at 100 iterations are those of the cells, so the comparison favours the published points. The count at 100 iterations moved between runs of the benchmark of the same release, from 3 to 4 to 0 in the record of the engine, which is the reading that this paper offers for a margin that is at most \benchWorstAbsHundred{} in absolute value in the file read here. Fourth, the knee rule of Eq.~\ref{eq:knee} returns the uniform end in \kneeZeroSeeds{} of \nSeeds{} seeds, so the single hologram that the engine returns is the
refined start and the front is the product of the method; and the target-spaced schedule, which the theory suggests, loses \ablTargetDiff{} of hypervolume because it stops short of the efficient tail.

Two results are positive but narrower than they look. The hypervolume gain of \hvDiff{} is the same in the three reference points tested in absolute terms and is positive in every seed, but it is computed against four points of the published methods and their one-parameter fronts, not against the best attainable set, which the front does not reach. And the time of the front, \timeAloneMin{} to \timeAloneMax{} seconds alone on this machine, is under the 60 seconds of the acceptance criterion with a thin margin, while a direct search sweep takes \dsSec{} seconds and returns one point.

\subsection{What the results imply for use}\label{sec:disc-use}

For an experimenter who needs a hologram of a trap array, the method gives a picture that the single-point methods do not: the cost in light of each step of uniformity, and the position of the published points relative to it. On the ten by ten task the picture says that weighted GS or phase-fixed WGS is the right choice if uniformity matters at all, because the extra \eMaxRatioPct{} percent of light costs almost all of the uniformity, and that GS and the adaptive additive iteration are inside the front. The front is worth computing when a decision depends on how much uniformity can be given up for light: when the power budget is tight, when the traps are shallow and the atoms
are loaded only if the depth is sufficient, or when an array larger than the task changes the position of the weighted GS point; the paper does not measure those cases, and the benchmark shows that on all six of its geometries, the ten by ten lattice included, the stored front is not above the published union at 100 iterations.

For the studio and for the method, the results suggest four changes that the paper does not implement. The start of the front should be the converged hologram of weighted GS and not 30 iterations of phase-fixed WGS, which removes the dominated points. The knee rule should not be the product $eu$ but a point that the user selects, or the largest efficiency subject to a uniformity floor. The converged row of the benchmark should give the front the weights it needs, since a stored front that ends at a weight of 0.05 cannot be compared with a union of points. And a device model should be put in the loss, as the crosstalk
and quantization models of the literature~\cite{schroff,persson} allow, because every number of this paper is for the ideal device.
